\MajorSection{القسم السادس}

\Couplet{١٤٥}{كل إيمان باطل، وكل إيمان حق؛ والحقيقة مرآة محطمة مبعثرة}{في شظايا لا تُحصى، فيما يعتقد كل واحد أن شظيته الصغيرة تحوي الكل.}{}

\Couplet{١٤٦}{ما الحقيقة؟ سؤال طُرح قديمًا. الجواب: الحقيقة الموضوعية كلها واحدة؛}{كما يصنع النصفان دائمًا كلًا؛ أما حقيقة أخلاقية للجميع فلا وجود لها.}{}

\Couplet{١٤٧}{يا قليلي العلم من الزهّاد، تعلّموا من أفلاطون وأرسطو؛}{فكما أن الحق واقعي كخيركم، فإن غير الحق واقعي أيضًا كالشر.}{\BurtonNote{حاشية برتون: \textenglish{Aflatûn} و\textenglish{Aristû} هما أفلاطون وأرسطو.}}

\Couplet{١٤٨}{كقصر تنعكس صورته في المجرى، وكبخار يمتزج بالسماء؛}{هكذا ينسج دماغ الإنسان الفاني الشبكة المتشابكة من الحق والأكاذيب.}{}

\Couplet{١٤٩}{ماذا نرى هنا؟ صورًا، لا أكثر! الصور تملأ أضوأ العيون وأقواها؛}{ولا نعرف الجوهر؛ وبين الظلال نعيش ونموت، ونحن أنفسنا ظلال.}{}

\Couplet{١٥٠}{أسمع: «الإيمان ينقل الجبال»؛ وأرى عمل العالم لا يبالي}{بتلك المفاخرة الحمقاء والتبجح الصاخب، اللذين يغذيان غرورنا.}{}

\Couplet{١٥١}{«الإيمان ثابت لا يتحرك»؛ ولماذا؟ لأن خيالات الإنسان السخيفة لا تزال باقية؛}{وستبقى إلى أن يحتقر الإنسان الأحكم أحلام يقظة شبابه.}{}

\Couplet{١٥٢}{تقولون: «مبارك أن يؤمن المرء»؛ وقد يكون في القول من الصدق ما يكفي؛}{وقد يضيف نورًا إلى الحياة؛ ولا يبقى إلا أن تبيّنوا لنا كيف.}{}

\Couplet{١٥٣}{حتى لو استطعت لما آمنت بحكاياتكم وخرافاتكم البالية المبتذلة؛}{إنها مضجرة كلحن أُعيد غناؤه، فيتعب الأذن الكليلة لامرئ ناعس.}{}

\Couplet{١٥٤}{حرية إرادة الإنسان مع علم الله السابق! أي نمو مسخي في الدماغ البشري؛}{وأي قوى من نور تستطيع اختراق هذا اللغز المتكاثف بالكلمات الخاوية؟}{}

\Couplet{١٥٥}{عبثًا يستنجد القلب بالعناية الإلهية؛ وليس من الحكمة طلب عون كهذا؛}{فعلى الإنسان أن يقر بالناموس العديم الرحمة الذي يسوس الكرة الأرضية والسماوات السبع.}{}

\Couplet{١٥٦}{«كونوا صبيانًا صالحين، واطلبوا الجنة، وتعالوا ادفعوا للكاهن الذي يحمل المفتاح»؛}{هكذا قال، ويقول، وسيظل يقول آخر من يدخل الجنة: هو.}{}

\Couplet{١٥٧}{أهذه كلمات جديرة بأن يسمعها الرجال؟ ومع ذلك فهذا لسان الكنيسة المعتاد؛}{صيحة العلق النهمة، قوية عالية، بين مزاميرها وأناشيدها الصاعدة إلى السماء.}{}

\Couplet{١٥٨}{ماذا؟ أيكون الإيمان فضلًا وحقًا يُطالب به، وهو يولد وينمو مع الدماغ؛}{امضِ، أيها الأحمق، في طريقك الأحمق، واغمس الموتى المدفونين في الماء المقدس!}{}

\Couplet{١٥٩}{لكن لا تتّبع طريق انعدام الحكمة، ولا تتشبث بهذا وتتبرأ من ذاك؛}{آمن بكل ما يؤمن به الإنسان؛ فالكل والعدم هنا سواء.}{}

\Couplet{١٦٠}{ولكن أحقًا يكون الأمر كذلك؟ وكيف لنا أن نعرف؟ لعل هذا القدر، هذا الناموس؛}{ليس إلا كلمة أو صوتًا أو نفَسًا؛ أو، في أحسن الأحوال، نظرية الزاهد المصاب بهوس القمر.}{}

\Couplet{١٦١}{نعم، قد تكون الحقيقة موجودة، ولكنها ليست هنا؛ وعلى البشر أن يطلبوها ويجدوها هناك؛}{لكن أين؟ لا أنا ولا أنت نستطيع أن نقول، ولا أي شيء ولدته الأم الأرض.}{}

\Couplet{١٦٢}{يكفي أن نفكر أن الحقيقة يمكن أن توجد؛ تعالَ نجلس حيث تتوهج الورود؛}{فإنه لا يعرف كيف يعرف، حقًا، من لا يعرف أيضًا كيف يتخلّص مما يعرف.}{}
